\chapter{Common health questions}
\section{Learning objectives}
Apply risk--benefit reasoning to common supplement questions; communicate uncertainty without dismissing symptoms; and identify when a question requires urgent or specialist assessment.

\section{Can vitamins provide energy?}
Calories provide energy. B vitamins help release and use that energy, but taking extra B vitamins does not act like a stimulant when intake is already adequate. Persistent fatigue warrants assessment of sleep, mood, anemia, thyroid disease, infection, medication effects, and other causes.

\section{Do supplements prevent chronic disease?}
The strength of evidence depends on the population, dose, comparator, duration, and outcome. A statistically significant result may be too small to matter clinically, and evidence from a food pattern should not automatically be transferred to an isolated pill.
Some targeted supplements prevent specific deficiency-related outcomes, such as folic acid reducing neural-tube-defect risk before and during early pregnancy. For broad prevention, evidence is mixed and often does not support high-dose multivitamins or antioxidant products as substitutes for food, exercise, sleep, vaccination, and evidence-based medical care. A supplement claim should name a dose, population, outcome, and duration.

\section{Is a multivitamin necessary?}
It can serve as nutritional insurance for some people with low intake, poor appetite, or a restricted diet, but it is not required for everyone. Compare its doses with other products and avoid formulations that supply several times the Daily Value without a clear reason.

\section{What should I do with a low test result?}
Confirm what the test measures, whether the sample and units are appropriate, and whether inflammation or another illness affects interpretation. Ask what dose, duration, food changes, and follow-up are appropriate. Do not continue a high dose indefinitely simply because a number was once low.

\section{The central message}
Health comes from systems: adequate food, safe water, movement, sleep, relationships, preventive care, and treatment of disease. Vitamins and minerals are essential parts of that system, not shortcuts around it.
